% ============ 第七节 募集资金运用与未来发展规划 ==============================
% 募集资金运用概况与专户管理 → 各募投项目（备案/投资构成/效益）
% → 募投项目与现有业务的关系 → 未来发展规划。金额与第二节概览勾稽。
\gssection{募集资金运用与未来发展规划}

\subsection{募集资金运用概况}

经公司2026年第一次临时股东会审议通过，本次发行募集资金扣除发行费用后，将按
轻重缓急顺序投资于以下项目：

\biaodw{万元}
\begin{gstab}
\begin{longtable}{|C{10mm}|L{40mm}|L{26mm}|R{27mm}|R{27mm}|}
\hline
\thA{序号} & \thB{项目名称} & \thB{实施主体} & \thB{投资总额} & \thBw{27mm}{拟使用\newline 募集资金} \\ \hline
\endfirsthead
\hline
\thA{序号} & \thB{项目名称} & \thB{实施主体} & \thB{投资总额} & \thBw{27mm}{拟使用\newline 募集资金} \\ \hline
\endhead
1 & 新一代大模型研发及训练平台建设项目 & 发行人、示例智算 & 82,456.00 & 82,456.00 \\ \hline
2 & 智能体应用与行业解决方案研发项目 & 发行人 & 35,208.00 & 35,208.00 \\ \hline
3 & 补充流动资金 & 发行人 & 32,336.00 & 32,336.00 \\ \hline
\multicolumn{3}{|c|}{\bfseries 合计} & {\bfseries 150,000.00} & {\bfseries 150,000.00} \\ \hline
\end{longtable}
\end{gstab}

若本次发行实际募集资金净额少于拟投入募集资金总额，不足部分由公司通过自有资
金或自筹方式解决；若实际募集资金净额超过拟投入总额，超出部分将按照中国证监
会及上交所关于募集资金管理的有关规定使用。在本次发行募集资金到位前，公司可
根据项目进度以自筹资金先行投入，募集资金到位后予以置换。

公司已制定《募集资金管理制度》，本次发行募集资金将存放于董事会批准设立的专
项账户集中管理、专款专用，并由公司、保荐人与存放募集资金的商业银行签订三方
监管协议。

\subsection{募集资金投资项目具体情况}

\subsubsection{新一代大模型研发及训练平台建设项目}

本项目由发行人及全资子公司示例智算共同实施，总投资82,456.00万元，其中建设
投资68,240.00万元（以智算集群及配套软硬件购置为主）、铺底流动资金14,216.00
万元，建设期2年。项目已取得项目备案（示发改备〔2025〕第××号），并已按规
定办理环境影响登记备案。

本项目将建设自主可控的智算集群、模型训练平台与数据平台，支撑新一代基础模型
与多模态模型的持续迭代，并提升开放平台的推理服务承载能力，缓解公司算力资源
对外部租赁的依赖。经测算，项目达产后预计年均新增收入约96,000.00万元，税后
内部收益率约18.65\%，具有良好的经济效益。

本项目预计建设期为2年，项目实施进度安排具体如下：

\begin{gstab}
{\centering
\begin{tabular}{|L{34mm}|C{10mm}|C{10mm}|C{10mm}|C{10mm}|C{10mm}|C{10mm}|C{10mm}|C{10mm}|}
\hline
\thA{时间单位：月} & \thB{T+3} & \thB{T+6} & \thB{T+9} & \thB{T+12} & \thB{T+15} & \thB{T+18} & \thB{T+21} & \thB{T+24} \\ \hline
机房及基础设施建设 & \gt & \gt & \gt & & & & & \\ \hline
算力设备采购与部署 & & \gt & \gt & \gt & \gt & & & \\ \hline
训练与数据平台研发 & \gt & \gt & \gt & \gt & \gt & \gt & & \\ \hline
基础模型训练与调优 & & & \gt & \gt & \gt & \gt & \gt & \gt \\ \hline
平台联调与上线运行 & & & & & & \gt & \gt & \gt \\ \hline
\end{tabular}\par
\vspace{0.15\baselineskip}
{\zihao{-5}注：T代表项目开始时点。\par}}
\end{gstab}

\subsubsection{智能体应用与行业解决方案研发项目}

本项目由发行人实施，总投资35,208.00万元，全部为研发投入及配套建设投资，建
设期3年。项目已取得项目备案（示发改备〔2025〕第××号）。

本项目将围绕行业智能体平台、企业级知识引擎及重点行业解决方案开展研发与工程
化建设，并引进高层次研发人才。项目不单独核算经济效益，但将增强公司行业落地
能力与持续创新能力，为公司中长期发展提供技术与产品储备。

\subsubsection{补充流动资金}

本项目拟使用募集资金32,336.00万元补充流动资金。公司业务规模持续扩张，解决
方案类项目制业务及算力资源采购对营运资金占用较大，补充流动资金将满足公司业
务增长带来的营运资金需求，降低财务风险，增强抗风险能力。

\subsection{募集资金投资项目与现有业务的关系}

本次募投项目均围绕公司主营业务展开：训练平台项目夯实模型研发与推理服务的算
力底座，智能体与解决方案研发项目提升行业落地能力，补充流动资金支持业务规模
扩张。公司在技术、人员、市场等方面已为募投项目的实施做好储备，不存在募投项
目与现有主营业务显著不同或跨界经营的情形。

\subsection{未来发展规划}

\subsubsection{总体发展战略}

公司将坚持``模型驱动、行业深耕''的发展战略，以基础模型能力为核心，持续完善
``模型研发—平台服务—行业落地''的业务体系，致力于成为国际领先的人工智能大
模型企业。

\subsubsection{具体发展计划}

\paragraph{技术研发计划}
持续加大基础模型、多模态与智能体方向的研发投入，依托募投项目完善训练与数据
基础设施，保持模型能力的行业领先性。

\paragraph{市场开拓计划}
深化金融、政务、制造等重点行业的标杆客户合作，扩大开放平台开发者生态，稳步
推进智能应用的用户增长与商业化。

\paragraph{基础设施建设计划}
按计划推进智算集群与训练平台建设，优化算力多源供应结构，提升单位算力的训练
与推理效率。

\paragraph{人才发展计划}
完善人才引进与培养体系，通过员工持股平台等长效激励机制吸引和保留核心人才，
支撑公司持续发展。

\paragraph{合规与安全治理计划}
持续完善数据合规、算法安全与内容安全治理体系，确保业务发展符合人工智能领域
法律法规及监管要求。

上述发展规划的实现存在不确定性，不构成公司的盈利预测或业绩承诺。
